\textbf{- Reference:} 
\href{https://e.math.cornell.edu/people/belk/numbertheory/CyclotomicPolynomials.pdf}{Fields and Cyclotomic Polynomials}~\cite{cyclotomic-polynomial}

\subsection{Definitions}
\label{subsec:cyclotomic-def}
\begin{tcolorbox}[title={\textbf{\tboxdef{\ref*{subsec:cyclotomic-def}} Cyclotomic Polynomial}}]
 \textbf{The $\bm{n}$-th Cyclotomic Polynomial:} is a polynomial whose roots are the primitive $n$-th roots of unity, that is: 
 
 \[ \Phi_n(x) = \prod_{\zeta \in P(n)} (x - \zeta)  = \prod_{\substack{0 \leq k \leq n-1,\\ 
 \text{gcd}(k, n) = 1}} (x - \omega^k) \textcolor{white}{......} \text{, where } \omega = e^{2\pi i/n} \]

 Remember the Euler's formula: $e^{2k\pi i/n} = \cos\left(\dfrac{2k\pi}{n}\right) + i\cdot\sin\left(\dfrac{2k\pi}{n}\right)$

$ $

A few pre-computed cyclotomic polynomials are as follows:
 
\begin{multicols}{2}
$\Phi_1(x) = x - 1$
\newline $\Phi_2(x) = x + 1$
\newline $\Phi_3(x) = x^2 + x + 1$
\newline $\Phi_4(x) = x^2 + 1$
\newline $\Phi_5(x) = x^4 + x^3 + x^2 + x + 1$
\newline $\Phi_6(x) = x^2 - x + 1$
\newline $\Phi_7(x) = x^6 + x^5 + x^4 + x^3 + x^2 + x + 1$
\newline $\Phi_8(x) = x^4 + 1$
\newline $\Phi_9(x) = x^6 + x^3 + 1$
\newline $\Phi_{10}(x) = x^4 - x^3 + x^2 - x + 1$
\end{multicols}
\end{tcolorbox}


As one example, 

\[\Phi_4(x) = \prod_{\substack{0 \leq k \leq 3,\\ 
 \text{gcd}(k, 4) = 1}} (x - \omega^k) = (x - \omega^1)(x - \omega^3) = (x - e^{2\pi i / 4})(x - e^{2\cdot3\pi i / 4}) = (x - e^{\pi i / 2})(x - e^{3\pi i / 2})\]

$= \left(x - \left(\cos\left(\dfrac{\pi}{2}\right) + i \cdot \sin\left(\dfrac{\pi}{2}\right) \right) \right) \cdot \left(x - \left(\cos\left(\dfrac{3\pi}{2}\right) + i \cdot \sin\left(\dfrac{3\pi}{2}\right) \right) \right)$

$= (x - i)(x + i) = x^2 + 1$



\subsection{Theorems}
\label{subsec:cyclotomic-theorem}


\begin{tcolorbox}[title={\textbf{\tboxtheorem{\ref*{subsec:cyclotomic-theorem}.1} Roots of the $M$-th Cyclotomic Polynomial}}]

Suppose that $M$ is a power of 2 and the $M$-th cyclotomic polynomial $\Phi_M(x) = x^n + 1$ (where $M = 2n$). Then, the roots of the $M$-th cyclotomic polynomial are $\omega, \omega^3, \omega^5, \cdots, \omega^{2n-1}$, where $\omega = e^{i\pi/n}$

\end{tcolorbox}

\begin{proof}

According to Definition~\ref*{subsec:cyclotomic-def} in \autoref{subsec:cyclotomic-def}, the roots of $\Phi_M(x)$ are $e^{2k\pi i/M} = e^{2k\pi i/(2n)} = e^{k\pi i/n}$ where $0 \leq k < M = 2n$ and $\textsf{gcd}(k, M = 2n) = 1$, thus $k = \{1, 3, 5, \cdots, 2n-1\}$. If we let $\omega = e^{i\pi/n}$, then the roots of $\Phi_M(x)$ are $\omega, \omega^3, \omega^5, \cdots, \omega^{2n-1}$.

\end{proof}

\begin{tcolorbox}[title={\textbf{\tboxtheorem{\ref*{subsec:cyclotomic-theorem}.2} Polynomial Decomposition into Cyclotomic Polynomials}}]
For any positive integer $n$, \[ x^n - 1 = \prod_{d \mid n} \Phi_d(x) \]
\end{tcolorbox}
\begin{myproof}
\begin{enumerate}
    \item The roots of $x^n - 1$ are all the $n$-th roots of unity. Thus, $x^n - 1 = (x - \omega^0)(x - \omega^1)...(x - \omega^{n-1})$, where $\zeta = \omega^k$.
    \item Theorem~\ref*{subsec:roots-theorem}.3 states that each $n$-th root of unity ($\omega^k$) is a primitive $d$-th root of unity for some $d$ that divides $n$. In other words, each $n$-th root of unity belongs to some $P(d)$ where $d \mid n$. Meanwhile, by definition, $\Phi_d(x) = \Pi_{\zeta \in P(d)} (x - \zeta)$. Therefore, $x^n - 1$ is the product of all $\Phi_d(x)$ such that $d \mid n$. 
\end{enumerate}
\end{myproof}
\begin{tcolorbox}[title={\textbf{\tboxtheorem{\ref*{subsec:cyclotomic-theorem}.3} Integer Coefficients of Cyclotomic Polynomials}}]
A cyclotomic polynomial has only integer coefficients.
\end{tcolorbox}
\begin{myproof}
We use strong induction on $n$. For $n=1$, $\Phi_1(x)=x-1$ is monic with integer coefficients. For $n>1$, assume the result for every positive $d<n$ and set
\[
F(x)=\prod_{\substack{d\mid n\\d<n}}\Phi_d(x).
\]
Then $F(x)$ is monic in $\mathbb{Z}[x]$. By Theorem~\ref*{subsec:cyclotomic-theorem}.2, $x^n-1=F(x)\Phi_n(x)$ in $\mathbb{C}[x]$. Long division by the monic polynomial $F$ uses only integer arithmetic, so its quotient and remainder lie in $\mathbb{Z}[x]$. Uniqueness of polynomial division over $\mathbb{C}$ shows that this quotient is $\Phi_n$ and the remainder is zero. Thus $\Phi_n\in\mathbb{Z}[x]$, completing the induction.
\end{myproof}
\begin{tcolorbox}[title={\textbf{\tboxtheorem{\ref*{subsec:cyclotomic-theorem}.4} Formula for $\bm{\Phi_{nk}(x)}$}}]
If $k \mid n$, then $\Phi_{nk}(x) = \Phi_n(x^k)$.
\end{tcolorbox}
\begin{myproof}
For positive integers $k\mid n$, Theorem~\ref*{subsec:order-theorem}.3 gives $\operatorname{ord}(\zeta)=nk$ if and only if $\operatorname{ord}(\zeta^k)=n$. Consequently, $\Phi_n(x^k)$ and $\Phi_{nk}(x)$ have exactly the same roots. These roots are simple: every root $\eta$ of $\Phi_n$ is nonzero and simple, and each polynomial $x^k-\eta$ has $k$ distinct complex roots because its derivative $kx^{k-1}$ is nonzero there. Both polynomials are monic and have the same simple roots, so they are equal. The map $\zeta\mapsto\zeta^k$ is $k$-to-one on these root sets, not one-to-one; multiplying once for every preimage would incorrectly repeat the roots of $\Phi_n$.
\end{myproof}